\section{Education}
\label{sec:education}

This section asks what the platform of \cref{sec:platform} could offer education, and where the
evidence stops. The rule from \cref{sec:platform-decisions} governs everything below: only a claim
labeled \code{observed_human_evidence}, or otherwise criterion-supported, may reach a learner. A
gap candidate --- always labeled \code{inferred} --- may motivate a question to an expert. It may
never become a fact taught to a student.

\subsection{Expert blind spots in education}
\label{sec:education-blindspots}

Three findings motivate the pipeline below, each given in full with its numbers in
\cref{sec:problem,sec:foundations}, so they are not repeated here. Physics experts categorize
problems by governing principle where novices use surface features
\parencite{chi1981categorization}, the reason the frozen Validation Track A study
(\cref{sec:validation}) chose principle selection as its topic. Physics teaching assistants (65\% of the maximum score) and instructors (68\%) predict students'
most common wrong answer well above the 40\% random-guessing baseline, with no significant
difference between the two groups \parencite{maries2016teaching}. This is why the pipeline below
routes ``what is difficult'' through learner data alone, never through expert or instructor
judgment. The
blind spot recurs across settings outside physics as well \parencite{nathan2003expert}. \evlit\ for
all three.

A worked example of tacit omission that the platform's design already takes seriously comes from
one university-physics intelligent tutoring system evaluation, with an hour-exam rubric fine enough
to see individual solution steps. Students who did their homework in Andes scored higher than
paper-homework students on the Drawings subscore (average effect size 1.21) and the Variable
definitions subscore (0.69), the two practices Andes requires, but not on Equations (0.11) or
Answers ($-0.08$) \parencite[Table~5]{vanlehn2005andes}.
\evlit. \evint\ The
steps instructors do not think to state explicitly are the steps a tutor that enforces
them improves --- and this fits the kind of step a methodology lens (\cref{sec:n3}) is designed to
flag as attested-but-unstated.

\subsection{The pipeline, and where its evidence runs out}
\label{sec:education-pipeline}

\Cref{tab:edu-pipeline} traces the full path from a gap map to a change in what a learner practices,
with the evidence status of each step stated plainly. The pipeline is evidence-complete only as far
as ``instructional material.'' Every step after that needs learner data the project does not have
(the NOW program runs with zero participants). Even with that data, every later step's own
literature shows a modest effect ceiling, not a large one.

\begin{table}[htbp]
  \centering
  \small
  \caption{The education pipeline: what each step does, and its evidence status.}
  \label{tab:edu-pipeline}
  \begin{tabularx}{\linewidth}{@{}p{0.16\linewidth}X@{}}
    \toprule
    Step & Evidence status \\
    \midrule
    Gap map (\cref{sec:n3}) &
      \evobs\ Runs end to end on three ledgers. \evint\ Prediction is not demonstrated: adversarial
      review found precision 3 of 19 strict on the pre-final maps, and closure does not beat a
      mismatched-evidence control on any ledger. \\
    Expert questions &
      \evobs\ The pipeline emits ranked, template-generated questions with a predicted knowledge
      type and elicitation channel. \evhyp\ That the questions target real residual is untested;
      the residual-measurement gate (\cref{sec:residual}) has not run. \\
    Captured residual &
      \evfut. No human answers exist outside the frozen Track A study, which is separately
      registered and does not depend on N3. \\
    Instructional material &
      \evlit\ Cognitive-task-analysis-based instruction outperforms other or conventional methods
      for identifying content, in a meta-analysis of a small number of studies (Hedges'
      $g = 0.871$) \parencite{tofelgrehl2013cognitive}. \evlit\ A worked example with a prompt that
      makes the learner generate a step beats one that states the step for them (self-explanation
      $g = 0.35$ over instructional explanation, $k = 6$) \parencite[Table~1]{bisra2018selfexplanation};
      simply adding more explanatory text to a worked example gives minimal benefit on its own
      \parencite{wittwer2010explanations}. \\
    Learner diagnostics &
      \evlit\ Must use learner data, not expert judgment or expert-authored text
      (\cref{sec:education-blindspots}). \evlit\ A single item response is a poor diagnostic unit:
      31\% of Force Concept Inventory answers changed on a one-week retest despite a total-score
      reliability of $r = 0.89$ \parencite{lasry2011fci}. \\
    Adaptive practice &
      \evlit\ Bayesian knowledge tracing models mastery from graded responses
      \parencite{corbett1995knowledge}; on next-response prediction, Bayesian knowledge tracing
      with extensions from the literature performed indistinguishably from deep knowledge tracing
      \parencite{khajah2016deep}, and logistic regression
      led on datasets of moderate size \parencite{gervet2020deep}. Learning gains from
      data-driven redesign are real but shown only on an immediate post-test in one selected unit
      \parencite[$d = 0.47$, $N = 91$]{liu2017closing}, and a 2026 study on units chosen by topic
      found no difference in learning gains ($N = 123$), though process measures favored the
      redesign \parencite{lyu2026redesign}. We found no study
      that compared a knowledge-tracing-driven policy against a simple mastery rule on a delayed or
      transfer outcome. \\
    \bottomrule
  \end{tabularx}
\end{table}

\subsection{Where established constructs enter}
\label{sec:education-constructs}

Pedagogical content knowledge, the form of subject knowledge specific to teaching it
\parencite{shulman1986those,shulman1987knowledge}, names the layer the gap map's L2 and L5 layers
target. This applies if the residual ever moves from an undocumented procedure to a teacher's
undocumented way of explaining it. \evint. Catalogued misconceptions
\parencite{smith1994misconceptions,disessa1993toward,vosniadou1992framework} are candidate content
for a lens's target construct, but the pipeline can only place them, not discover that a learner
holds one; that step is learner diagnostics, above, and is not built. Threshold concepts
\parencite{meyer2003threshold,meyer2005threshold} are a strong candidate description for a gap
map's highest-value item, if one is ever found; this report does not claim it has. Worked examples
and self-explanation prompts are the mechanism the pipeline currently targets, on the strength of a
broader literature \parencite{barbieri2023metaanalysis,sweller1988cognitive}. Intelligent tutoring
systems are both a comparison class ($d = 0.76$ over no tutoring, close to human tutoring's 0.79,
\parencite{vanlehn2011relative}; a randomized trial found an AI physics tutor outperformed in-class
active learning, \parencite{kestin2025aitutoring}) and a candidate delivery mechanism for
``adaptive practice,'' should it run. Knowledge tracing is the established technique for that step,
and its ceiling, not its existence, is the open question.

\subsection{Synthetic learners: a predictor, never a criterion}
\label{sec:education-synthetic}

A model prompted to answer as a student is not a substitute for a student's actual response, and
the physics-specific evidence for that claim is direct. A GPT-4 simulation of the Force Concept
Inventory showed \emph{no response variance at all} when told only to answer as a different
cohort. It produced variance that approximated human responses in some respects only once it was
told which preconception to hold \parencite{kieser2023chatgpt}. \evlit. A synthetic cohort therefore returns the misconception
taxonomy it was handed; it cannot discover a new one, and it cannot report how common one actually
is. The wider literature the reorientation reviewed agrees. Synthetic survey responses vary less than
real ones \parencite{bisbee2024synthetic}. Language-model replications of 156 scenario-based
psychology and management experiments gave larger effect sizes overall and significant results in
68--83\% of the cases whose original finding was null \parencite{cui2025large}. \evlit. Synthetic learners therefore belong only
at the instrument-debugging and candidate-hypothesis tiers: generating a probe to pilot before a
real student sees it, or checking that a pipeline runs at all on known-truth input. They are never
evidence for a gate, a power calculation, or a claim about how common a difficulty is. This is the
learner-layer restatement of the rule \cref{sec:platform-decisions} already states for the gap map
itself: a synthetic output may predict, and it may never decide.

\subsection{How Validation Track A fits, unchanged}
\label{sec:education-tracka}

The frozen Validation Track A study (\cref{sec:validation}) is not an output of the gap map and
does not depend on N3 succeeding. It is a separately registered study that starts only on its own
preconditions --- recruited experts, a course, ethics approval --- and its registered rules run
unchanged regardless of anything in this section. Stage A (twelve physicists, AI-led and
human-led retrospective probing, counterbalanced) is the human-answered version of ``expert
questions $\to$ captured residual'' in \cref{tab:edu-pipeline}; its pre-registered K0 check, that
at least 30\% of coded student errors are principle-selection or representation errors, is the
human-data confirmation that a candidate bottleneck is real. This is exactly the confirmation an
inferred gap candidate would need before anyone trusted it. Stage B, a two-arm randomized trial
with a delayed-transfer outcome, is the ``does the captured residual, taught, move a learning
outcome'' test. It runs as a pre-registered futility test at a smallest-effect-of-interest of
$d = 0.30$, because the worked-examples literature above predicts a true effect below that
threshold is plausible. Nothing in this report changes a single Track A number, and no addition
from the NOW program may gate Track A's own readings (\cref{sec:validation}). Any comparison
between a sealed, unopened gap-map reconstruction of Track A's topic and Stage A's independently
coded operations is registered as secondary and non-gating.

\subsection{A candidate education pilot}
\label{sec:education-pilot}

A small pilot could test whether a gap-map-derived instructional addition changes anything
measurable, at the smallest defensible scale, before any learner-facing deployment is proposed.
One course section cannot power a test against typical education effect sizes: across 1{,}942
effect sizes from 747 randomized trials with standardized achievement outcomes, the median was
0.10 SD \parencite{kraft2020interpreting}. This is a feasibility pilot only, not powered to detect an
effect, and any confidence interval it reports should be read as such. One introductory mechanics
course section would run two arms of the same worked-example module,
matched on word count and prompt count. The treatment arm's self-explanation prompts would carry
one or two operations that the gap map proposed \emph{and} that experts then confirmed in blind
elicitation (\cref{sec:validation}). Only operations with that human confirmation qualify;
unvalidated candidates stay out of learner material. The primary measure would be delayed transfer performance, scored by
instructors blind to which operations were added, guarding against test-alignment inflation
\parencite{kulik2016effectiveness}; logged time on task and prompt-completion rate would serve as a
manipulation check.

\textbf{What this pilot would not claim, on any result.}
\begin{itemize}[noitemsep]
  \item That the gap map found hidden expert knowledge on its own: expert confirmation shows only
    that experts hold this operation, not that the map finds such operations reliably elsewhere.
  \item Individual-learner diagnosis: the instrument is validated for cohort comparison, not one
    student's state, and single-item responses are unstable
    \parencite{lasry2011fci,madsen2017inventories}.
  \item A substitute for, comparison with, or gate on Validation Track A: a result here says
    nothing about whether AI-assisted expert elicitation works, only whether a text-derived
    candidate moves an outcome.
  \item A synthetic learner used beyond piloting a prompt.
\end{itemize}

The honored rule stands regardless of this pilot's result: an inferred gap candidate may become a
question worth asking an expert, and, once validated, an item worth teaching. It may not skip that
validation step and reach a learner directly, no matter how plausible the candidate looks or how
well the pilot's manipulation check passes.
